% Part: lambda-calculus
% Chapter: lambda-definablity
% Section: arithmetical-functions

\documentclass[../../../include/open-logic-section]{subfiles}

\begin{document}

\olfileid{lam}{rep}{arf}
\olsection{\usetoken{S}{lambda definable} अंकगणितीय फलने}

\begin{prop}
  \ollabel{prop:succ-ld}
  उत्तरवर्ती फलन~$\Succ$ हे !!{lambda definable} आहे.
\end{prop}

\begin{proof}
उत्तरवर्ती फलनाला !!{lambda define}s करणारे पद असे आहे:
\begin{align*}
  \fn{Succ} & \ident \lambd[a][\lambd[fx][f ({a} f x)]].
  \intertext{आपल्या संकेतनरूढीनुसार याचा पूर्ण अर्थ असा:}
  \fn{Succ} & \ident \lambd[a][\lambd[f][\lambd[x][(f (({a} f) x))]]].
  \intertext{$\fn{Succ}$ हे~$a$ हा संख्यारूपी फलसाधक
    स्वीकारणारे फलन आहे आणि त्याचे मूल्य दुसरे फलन,
    $\lambd[fx][f ({a} f x)]$, आहे. ते फलन स्वतः चर्च अंक
    नाही. मात्र $a$ हा फलसाधक चर्च अंक असेल, तर त्याचे
    न्यूनीकरण चर्च अंकात होते. पाहा:}
   (\lambd[a][\lambd[fx][f ({a} f x)]])\,\num{n} & \redone
   \lambd[fx][f ({\num{n}} f x)].
   \intertext{आतील पद $\num{n}fx$ मध्ये रेडेक्स आहे, कारण
     $\num{n}$ हे $\lambd[fx][f^nx]$ आहे. म्हणून
     $\num{n}fx \red f^nx$. त्यामुळे संपूर्ण पदासाठी}
   \fn{Succ}\, \num n & \red \lambd[fx][f(f^n(x))],
\end{align*}
म्हणजे $\num{n+1}$.
\end{proof}

\begin{ex}
  $\fn{Succ}$ हे~$\num{0}$, म्हणजे
  $\lambd[fx][x]$, ला लागू केल्यावर काय होते ते पाहू.
  पदे पूर्णपणे लिहू:
  \begin{align*}
    \fn{Succ}\, \num{0} & \ident (\lambd[a][\lambd[f][\lambd[x][(f (({a} f) x))]]])(\lambd[f][\lambd[x][x]])\\
    & \redone \lambd[f][\lambd[x][(f (({(\lambd[f][\lambd[x][x]])} f) x))]] \\
    & \redone \lambd[f][\lambd[x][(f ({(\lambd[x][x])} x))]] \\
    & \redone \lambd[f][\lambd[x][(f x)]] \ident \num{1}\\
  \end{align*}
\end{ex}

\begin{prob}
  पुढील पद
  \begin{align*}
    \fn{Succ}' & \ident \lambd[{n}][\lambd[fx][{n} f (f x)]]
  \end{align*}
  उत्तरवर्ती फलनाला !!{lambda define}s करते. का ते स्पष्ट करा.
\end{prob}

\begin{prop}
  \ollabel{prop:add-ld}
  बेरीज फलन~$\Add$ हे !!{lambda definable} आहे.
\end{prop}

\begin{proof}
बेरीज पुढील पदांनी !!{lambda define}d होते:
\begin{align*}
  \fn{Add} & \ident \lambd[{a}{b}][\lambd[fx][{a} f ({b} f x)]]
\intertext{किंवा पर्यायाने,}
  \fn{Add}' & \ident \lambd[{a}{b}][{a}\, \fn{Succ} \, {b}].
\intertext{पहिले बेरीजपद असे कार्य करते: $\fn{Add}$ प्रथम
  ${a}$ आणि ${b}$ या दोन संख्या स्वीकारते. त्यातून
  $f$ आणि $x$ स्वीकारून $af(bfx)$ परत करणारे फलन
  मिळते. $a$ आणि $b$ हे चर्च अंक $\num{n}$ आणि
  $\num{m}$ असतील, तर त्याचे न्यूनीकरण $f^{n+m}(x)$
  मध्ये होते; ते $f^{n}(f^{m}(x))$ सारखेच आहे.
  टप्प्याटप्प्याने पाहू:}
  (\lambd[{a}{b}][\lambd[fx][{a} f ({b} f x)]])\num n\,\num m & \red
  \lambd[fx][\num{n}\, f (\num {m}\, f x)] \\
  & \red \lambd[fx][\num{n}\, f (f^m x)] \\
  & \red \lambd[fx][f^n (f^m x)] \ident \num {n+m}.
\intertext{बेरीजचे दुसरे निरूपण $\fn{Add'}$ वेगळे कार्य
  करते: ते चर्च अंक $\num{n}$ आणि $\num{m}$ यांना
  लागू केल्यास,}
\fn{Add}' \num n \,\num m
& \red \num{n}\, \fn{Succ}\, \num{m}.
\intertext{पण $\num{n} f x$ चे नेहमी $f^n(x)$ मध्ये
  न्यूनीकरण होते. म्हणून,}
  \num{n}\,  \fn{Succ}\, \num{m} & \red \fn{Succ}^n(\num{m}).
\end{align*}
$\fn{Succ}$ हे उत्तरवर्ती फलनाला !!{lambda define}s
करते आणि~$m$ ला उत्तरवर्ती फलन $n$ वेळा लागू
केल्यावर $n+m$ मिळते. म्हणून याचे पुढे
$\num{n+m}$ मध्ये न्यूनीकरण होते.
\end{proof}

\begin{prop}
  \ollabel{prop:mult-ld}
  गुणाकार पुढील पदाने !!{lambda definable} आहे:
  \[
  \fn{Mult} \ident \lambd[ab][\lambd[fx][a (b f) x]]
  \]
\end{prop}

\begin{proof}
  हे कसे कार्य करते ते पाहण्यासाठी $\fn{Mult}$ ला
  चर्च अंक $\num{n}$ आणि $\num{m}$ लागू करू.
  $\fn{Mult} \, \num{n} \, \num{m}$ चे न्यूनीकरण
  $\lambd[fx][\num{n}(\num{m}\, f)x]$ मध्ये होते.
  पद $\num{m} f$ हे आपल्या फलसाधकाला $f$ हे
  $m$ वेळा लागू करणारे फलन परिभाषित करते.
  परिणामी, $\num{n} (\num{m} f) x$ हे
  ``$f$ हे $m$ वेळा लागू करा'' या फलनाला~$x$ वर
  $n$ वेळा लागू करते. म्हणजे $f$ हे $x$ वर
  $n\cdot m$ वेळा लागू होते. मिळणारे सामान्य पद
  म्हणजे नेमका चर्च अंक $\num{nm}$.
\end{proof}

\begin{editorial}
$\eta$-न्यूनीकरणाने हे पद आणखी सोपे करता येते:
\[
  \fn{Mult} \ident \lambd[ab][\lambd[f][a (b f)]].
\]

पण त्याआधी $\eta$-न्यूनीकरण स्पष्ट करावे लागेल.
\end{editorial}

\begin{prob}
पुढील पदाने गुणाकार !!{lambda define}d होतो:
\[
  \fn{Mult}' \ident \lambd[ab][b (\fn{Add}\, a) \num{0}].
\]
हे का कार्य करते ते स्पष्ट करा.
\end{prob}

%READERNOTE{SOL6-A128: मूळ e b हे पद शून्य घातांकासाठी identity मध्ये न्यूनीत होते; ते बीटा-सामान्य चर्च अंक एक नाही. बाहेरून चर्च अंकाचे दोन प्राचल लावल्याने सर्व घातांकांसाठी आवश्यक बीटा-सामान्य अंक मिळतो. शून्याच्या शून्य घाताचा येथे एक असा पुनरावर्ती संकेत आहे.}
घातांक फलनाची $\lambd$-पद म्हणून व्याख्या
आश्चर्यकारकपणे सोपी आहे:
\[
  \fn{Exp} \ident \lambd[be][\lambd[fx][e b f x]].
\]
पहिला फलसाधक $b$ हा पाया आणि दुसरा~$e$ हा
घातांक आहे. आपल्या संख्या-संकेतनानुसार $e f$
हे सहजपणे $f^e$ आहे. मात्र चर्च अंकाचे पूर्ण
बीटा-सामान्य रूप मिळण्यासाठी दोन बाह्य प्राचल आवश्यक
आहेत. घातांक शून्य असल्यास आतले पद identity मध्ये
न्यूनीत होते आणि या दोन प्राचलांनी चर्च अंक एक मिळतो.
इतर घातांकांसाठी पद वारंवार फलनसंयोजन करून योग्य
चर्च अंक देते. येथे शून्याच्या शून्य घाताचे मूल्यही एक
मानले आहे. हे समजणे कठीण वाटल्यास,
वारंवार गुणाकारानेही घातांक फलन परिभाषित करता येते:
\[
  \fn{Exp}' \ident \lambd[be][e (\fn{Mult}\, b) \num{1}].
\]

चर्च अंकांवरील पूर्ववर्ती फलन आणि वजाबाकी वाटतात
तितकी सोपी नाहीत: त्यांसाठी जोड्यांचे संकेतन आवश्यक आहे.

\end{document}
